\section{Appendix}

\subsection{Authorship and Contributions}

This report is partitioned so each major part has a primary author responsible for writing and technical correctness.

\subsubsection*{Deniz Isikli (s215818)}
Primary author of:
\begin{itemize}
    \item Chapter 3 (Design Choices and Information-Flow Argument),
    \item Chapter 5 (Security Analysis of API Commands),
\end{itemize}

\subsubsection*{Benas Skripkiunas (s253140)}
Primary author of:
\begin{itemize}
    \item Chapter 2 (Participants, Roles, and Assumptions),
    \item Chapter 6 (Declassification Justification),
\end{itemize}

\subsubsection*{Abid Hasan (s252729)}
Primary author of:
\begin{itemize}
    \item Chapter 1 (Introduction),
    \item Chapter 4 (Server API and Security Semantics),
\end{itemize}

\subsection{Appendix A: Flow Checker Pseudocode}

The server's type checker implements the following core operations, referenced in Chapters 3, 4, and 5.

\subsubsection*{Flow Order Checking}

The function \texttt{canFlow} checks whether information is allowed to flow from label $L_1$ to label $L_2$:

\begin{verbatim}
canFlow(label L1=(C1,I1), label L2=(C2,I2)):
  // Check confidentiality: for all (o:r1) in C1,
  // must exist (o:r2) in C2 with r1 ⊇ r2
  for each policy (o:r1) in C1:
    found = false
    for each policy (o:r2) in C2:
      if r1 ⊇ r2:
        found = true; break
    if not found:
      return false
  
  // Check integrity: for all (o <- w2) in I2,
  // must exist (o <- w1) in I1 with w1 ⊇ w2
  for each policy (o <- w2) in I2:
    found = false
    for each policy (o <- w1) in I1:
      if w1 ⊇ w2:
        found = true; break
    if not found:
      return false
  
  return true
\end{verbatim}

\subsubsection*{Label Join (Least Upper Bound)}

The function \texttt{join} computes the least upper bound of two labels (used when merging information from multiple sources):

\begin{verbatim}
join(label L1=(C1,I1), label L2=(C2,I2)):
  // Confidentiality join: union of all (o:r) from both
  // but when two policies for same owner exist,
  // keep only the stricter (smaller reader set)
  C_result = {}
  for each owner o:
    if (o:r1) in C1 and (o:r2) in C2:
      C_result += (o:(r1 ∩ r2))  // intersection for conf
    else if (o:r1) in C1:
      C_result += (o:r1)
    else if (o:r2) in C2:
      C_result += (o:r2)
  
  // Integrity join: union all (o <- w) from both
  // but when two policies for same owner exist,
  // keep only the weaker (larger writer set)
  I_result = {}
  for each owner o:
    if (o <- w1) in I1 and (o <- w2) in I2:
      I_result += (o <- (w1 ∪ w2))  // union for integ
    else if (o <- w1) in I1:
      I_result += (o <- w1)
    else if (o <- w2) in I2:
      I_result += (o <- w2)
  
  return (C_result, I_result)
\end{verbatim}

\subsubsection*{Implicit Flow Handling}

When type-checking a branch, the control-context label (program counter label $pc$) is updated:

\begin{verbatim}
check_assignment(env Gamma, label pc, 
                 var x, expr e):
  L_e = typeOf(Gamma, e)
  L_target = Gamma(x)
  
  // Assignment allowed iff (pc ⊔ L_e) ⊆ L_target
  if canFlow(join(pc, L_e), L_target):
    return OK
  else:
    return ERROR("Flow violation on line ...; 
                  type of e is too high, 
                  x cannot hold it")

check_if_branch(env Gamma, label pc,
                expr guard, block B):
  L_guard = typeOf(Gamma, guard)
  pc_new = join(pc, L_guard)
  
  // Check body with raised pc
  check_block(Gamma, pc_new, B)
\end{verbatim}

\subsection{Appendix B: Label Comparison Examples}

Below are five concrete examples of flow checking, worked through by hand:

\textbf{Example 1: Simple Confidentiality Flow}

Source: $L_1 = (doctor:\{doctor,nurse\})$
Destination: $L_2 = (doctor:\{nurse\})$

Is $L_1 \sqsubseteq L_2$?

\begin{itemize}
    \item Confidentiality: $(doctor:\{doctor,nurse\})$ in $L_1$; $(doctor:\{nurse\})$ in $L_2$. Check: $\{doctor, nurse\} \supseteq \{nurse\}$? Yes.
    \item Integrity: neither label has integrity policies.
\end{itemize}

Result: \textbf{Yes}, $L_1 \sqsubseteq L_2$. Information can flow.

\textbf{Example 2: Integrity Constraint Blocks Assignment}

Source: $L_1 = (hospital \leftarrow \{nurse\})$ (data only trusted from nurses)
Destination: $L_2 = (hospital \leftarrow \{doctor\})$ (data must be trusted from doctors)

Is $L_1 \sqsubseteq L_2$?

\begin{itemize}
    \item Confidentiality: neither has confidentiality policies.
    \item Integrity: $L_2$ requires $(hospital \leftarrow \{doctor\})$. In $L_1$, we have $(hospital \leftarrow \{nurse\})$. Check: $\{nurse\} \supseteq \{doctor\}$? No.
\end{itemize}

Result: \textbf{No}, $L_1 \not\sqsubseteq L_2$. A nurse cannot modify data that must come from doctors only.

\textbf{Example 3: Multi-Owner Confidentiality}

Source: $L_1 = (patient:\{patient,doctor\}) \wedge (hospital:\{hospital,doctor\})$
Destination: $L_2 = (patient:\{doctor\}) \wedge (hospital:\{hospital,doctor\})$

Is $L_1 \sqsubseteq L_2$?

\begin{itemize}
    \item Patient policy: $(patient:\{patient,doctor\})$ in $L_1$; $(patient:\{doctor\})$ in $L_2$. Check: $\{patient,doctor\} \supseteq \{doctor\}$? Yes.
    \item Hospital policy: $(hospital:\{hospital,doctor\})$ in $L_1$; $(hospital:\{hospital,doctor\})$ in $L_2$. Check: $\{hospital,doctor\} \supseteq \{hospital,doctor\}$? Yes.
\end{itemize}

Result: \textbf{Yes}, $L_1 \sqsubseteq L_2$. Data can flow because both owner policies satisfy the confidentiality order.

\textbf{Example 4: Implicit Flow via PC Raise}

Program counter: $pc = (secret:\{doctor\})$ (inside a branch guarded by secret data)
Source: $L_e = \emptyset$ (a constant, no label)
Target: $\Gamma(public) = (public:\{\})$ (public writable by anyone)

Check: $pc \sqcup L_e \sqsubseteq L_{target}$?

\begin{itemize}
    \item $pc \sqcup L_e = (secret:\{doctor\})$ (the join is dominated by $pc$).
    \item Is $(secret:\{doctor\}) \sqsubseteq (public:\{\})$? This requires: for $(secret:\{doctor\})$, there must be a policy $(secret:r')$ in the target with $\{doctor\} \supseteq r'$. But $L_{target}$ has no secret policy. 
    \item The check fails.
\end{itemize}

Result: \textbf{No}, assignment is rejected. Implicit flow from secret to public is blocked.

\textbf{Example 5: Join Operation}

Label 1 (input A): $L_A = (owner:\{owner,user1\})$
Label 2 (input B): $L_B = (owner:\{owner,user2\})$

Compute $L_A \sqcup L_B$ (the join; output label must accommodate both inputs):

\begin{itemize}
    \item For owner in confidentiality: $r_A = \{owner, user1\}$ and $r_B = \{owner, user2\}$.
    \item Join of readers: $r_A \cap r_B = \{owner\}$ (intersection for confidentiality; stricter is safer).
    \item Result: $L_A \sqcup L_B = (owner:\{owner\})$.
\end{itemize}

Interpretation: An output that depends on both $A$ and $B$ cannot be labeled more permissively than both; the result has the tighter (more confidential) label.

\subsection{Appendix C: Formal Myers-Liskov Label Semantics}

This section restates the formal definitions underlying the system, referenced in Chapter 2.

\subsubsection*{Label Semantics}

Let $\mathcal{A}$ be a fixed set of principals (actors). A \emph{confidentiality policy} has the form $(o:r)$ where $o \in \mathcal{A}$ is an owner and $r \subseteq \mathcal{A}$ is a reader set. Intuitively, $(o:r)$ means owner $o$ declares that only principals in $r$ are permitted to read this data.

A \emph{confidentiality label} is a conjunction of policies: $C = \bigwedge_i (o_i:r_i)$. For a principal $p$ to be authorized to read data with label $C$, $p$ must be in $r_i$ for all $i$. That is, $p$ must satisfy the constraints of all owners.

An \emph{integrity policy} has the form $(o \leftarrow w)$ where $o \in \mathcal{A}$ is an owner and $w \subseteq \mathcal{A}$ is a writer set. Intuitively, $(o \leftarrow w)$ means owner $o$ trusts only principals in $w$ to modify this data.

An \emph{integrity label} is a conjunction: $I = \bigwedge_j (o_j \leftarrow w_j)$. For a principal $p$ to be authorized to write, $p$ must be in $w_j$ for all $j$.

A \emph{security label} is a pair $L = (C, I)$ of a confidentiality label and an integrity label.

\subsubsection*{Information Flow Order}

The information-flow order is defined as:

\[
L_1=(C_1,I_1) \sqsubseteq L_2=(C_2,I_2) \quad \Leftrightarrow
\]

\begin{enumerate}
\item For each policy $(o:r_1) \in C_1$, there exists a policy $(o:r_2) \in C_2$ such that $r_1 \supseteq r_2$ (confidentiality can only become stricter).

\item For each policy $(o \leftarrow w_2) \in I_2$, there exists a policy $(o \leftarrow w_1) \in I_1$ such that $w_1 \supseteq w_2$ (integrity can only become more stringent).
\end{enumerate}

The relation is a partial order (reflexive, transitive, antisymmetric).

\subsubsection*{Least Upper Bound (Join)}

The join operation computes the least upper bound under $\sqsubseteq$:

\[
L_1 \sqcup L_2 = (C_1 \sqcup C_2, I_1 \sqcup I_2)
\]

where for confidentiality, $(o:r)$ in the result is $(o:(r_1 \cap r_2))$ if both labels have a policy for $o$, or the individual policy if only one does. The reader set shrinks (becomes stricter).

For integrity, $(o \leftarrow w)$ in the result is $(o \leftarrow (w_1 \cup w_2))$ if both labels have a policy for $o$, or the individual policy if only one does. The writer set expands (becomes more permissive, trusting more sources).

\subsection{Appendix D: API Trace Example}

Below is a worked example of a hospital scenario, showing how data flows through API commands while maintaining invariants.

\textbf{Scenario:} A patient uploads a medical image, a doctor encrypts it using a program, and then the hospital declassifies a summary for research.

\textbf{Step 1: UPLOAD (Patient uploads image)}

\begin{verbatim}
Input:
  token = <valid auth for patient>
  payload = <medical image bytes>
  C = (patient:{patient, doctor})
  I = (patient <- {patient})

Server checks:
  - token is valid ✓
  - C and I parse correctly ✓
  
Server effect:
  creates object:
    dataId = 1001
    label = ((patient:{patient,doctor}), (patient <- {patient}))
    owner = patient
    payload = <image>
    
Response:
  ok(1001)
\end{verbatim}

\textbf{Step 2: UPLOAD\_PROGRAM (Doctor uploads encryptor program)}

\begin{verbatim}
Input:
  token = <valid auth for doctor>
  src = "
    input image: (patient:{patient,doctor});
    output: (patient:{patient,doctor}) wedge (hospital:{hospital});
    
    encout := image;  // copy
  "
  signature = "doctor@hospital.org"

Server checks:
  - source parses ✓
  - input and output labels are valid ✓
  - Flow check: (patient:{patient,doctor}) ⊆ (patient:{patient,doctor}) ∧ 
    (hospital:{hospital})? No! (patient:{patient,doctor}) does not 
    include (hospital:{hospital}) policy.
  
Server response:
  error(flow_violation, "output label too high; 
  source has no hospital policy")
\end{verbatim}

The program is rejected. Doctor corrects it:

\begin{verbatim}
Input (corrected):
  src = "
    input image: (patient:{patient,doctor});
    output: (patient:{patient,doctor});
    
    encout := image;
  "

Server checks:
  - Flow: (patient:{patient,doctor}) ⊆ (patient:{patient,doctor})? Yes ✓
  
Server effect:
  creates program:
    programId = 2001
    interface = {input: (patient:{patient,doctor}), 
                 output: (patient:{patient,doctor})}
    
Response:
  ok(2001, {input: ..., output: ...})
\end{verbatim}

\textbf{Step 3: RUN\_PROGRAM (Doctor runs encryptor)}

\begin{verbatim}
Input:
  token = <valid auth for doctor>
  programId = 2001
  inputIds = [1001]  // the image from step 1
  outC = (patient:{patient, doctor})
  outI = (patient <- {patient, doctor})

Server checks:
  - programId 2001 exists ✓
  - dataId 1001 exists ✓
  - doctor can read 1001? (patient:{patient,doctor}), doctor in {patient,doctor}? Yes ✓
  - program interface: input requires (patient:{patient,doctor}) ✓
  - input labels: L_1001 = ((patient:{patient,doctor}), (patient <- {patient}))
  - required output lower bound: L_req = L_1001 = ((patient:{patient,doctor}), 
    (patient <- {patient}))
  - Is L_req ⊆ ((patient:{patient,doctor}), (patient <- {patient,doctor}))?
    Conf: (patient:{patient,doctor}) ⊆ (patient:{patient,doctor})? Yes ✓
    Integ: for (patient <- {patient,doctor}) in output, must exist 
           (patient <- w) in input with w ⊇ {patient,doctor}.
           Input has (patient <- {patient}). Is {patient} ⊇ {patient,doctor}? 
           No. ✗
    
Server response:
  error(flow_violation, "output integrity too high; 
  input is not trusted by enough writers")
\end{verbatim}

The doctor corrects the output label:

\begin{verbatim}
Input (corrected):
  outI = (patient <- {patient})

Server checks:
  - Is L_req ⊆ output? Yes ✓
  
Server effect:
  executes program in sandbox
  result bytes = <encrypted image>
  creates output object:
    dataId = 1002
    label = ((patient:{patient,doctor}), (patient <- {patient}))
    
Response:
  ok(1002)
\end{verbatim}

\textbf{Step 4: DECLASSIFY\_AND\_RELEASE (Hospital releases summary)}

\begin{verbatim}
Input:
  token = <valid auth for admin>
  dataId = 1002
  transform = "aggregate"
  params = {k = 5}
  targetC = (patient:{patient,doctor,researcher}) 
            wedge (hospital:{hospital,researcher})

Server checks:
  - 1002 exists, label is ((patient:{patient,doctor}), 
    (patient <- {patient})) ✓
  - transform "aggregate" is approved ✓
  - params k=5 satisfies k_min=5? Yes ✓
  - authority check:
    * for (patient:{patient,doctor}) → (patient:{patient,doctor,researcher}),
      admin must have authority from patient. Does admin have delegation?
      (not shown, assume yes in this example) ✓
    * for (hospital:{hospital}) → (hospital:{hospital,researcher}),
      admin must have authority from hospital. (admin typically has this) ✓
  
Server effect:
  applies aggregate transform:
  result = <anonymized aggregate summary>
  creates output object:
    dataId = 1003
    label = ((patient:{patient,doctor,researcher}), 
             (hospital:{hospital,researcher}))
  writes audit entry:
    (actor: admin, time: 2024-04-22 10:30, 
     source: 1002, target: 1003, 
     transform: aggregate, k: 5,
     authority: patient-delegation-token-xyz)
     
Response:
  ok(1003)
\end{verbatim}

\textbf{Invariant Verification:}

After all four commands:
\begin{itemize}
    \item \textbf{I1 (Label Validity)}: Objects 1001, 1002, 1003 all have valid labels. ✓
    \item \textbf{I2 (Read Safety)}: Only patients and doctors can read 1001/1002. Only patients/doctors/researchers can read 1003. Read checks enforce this. ✓
    \item \textbf{I3 (Write Safety)}: Only patients can write to 1001/1002. Integrity labels are respected. ✓
    \item \textbf{I4 (Program Safety)}: Program 2001 was flow-checked before upload; RUN\_PROGRAM verified label compatibility. ✓
    \item \textbf{I5 (Declassification Safety)}: Declassification in step 4 was authority-checked and audited. ✓
\end{itemize}
